\chapter*{Leitidee}
\addcontentsline{toc}{chapter}{Leitidee}

\textbf{Leitidee}

In \textbf{Mathematik 1} standen lokale Eigenschaften von Funktionen im Mittelpunkt.
Ableitungen beschreiben das Verhalten einer Funktion in der unmittelbaren Umgebung
eines Punktes. Damit lassen sich Tangenten, Extremstellen und Wendepunkte bestimmen
und lokale Dynamik analysieren.

Dar�ber hinaus haben wir grundlegende Bausteine kennengelernt:
Exponentialfunktionen modellieren Wachstum und Zerfall, und komplexe Zahlen
erm�glichen es, Drehungen und Schwingungen algebraisch zu beschreiben.
Mit der Eulerschen Identit�t
\[
e^{ix}=\cos(x)+i\sin(x)
\]
wird sichtbar, dass Exponentialfunktionen und harmonische Schwingungen
eng miteinander verbunden sind.

In \textbf{Mathematik 2} erweitern wir diese Perspektive.
Funktionen werden nicht mehr nur punktweise untersucht,
sondern als Elemente strukturierter R�ume betrachtet.
Damit wird es m�glich, Funktionen und Signale global zu beschreiben,
zu rekonstruieren und in einfachere Bausteine zu zerlegen.

Die zentralen Werkzeuge daf�r sind
\emph{Integration}, \emph{Projektionen in Funktionenr�umen} (Bestapproximation)
sowie die \emph{Fourieranalyse}.
Mit Hilfe geeigneter Transformationen lassen sich au�erdem dynamische Prozesse
und technische Systeme systematisch analysieren.

Der konzeptionelle �bergang l�sst sich so zusammenfassen:

\medskip

\begin{center}
\textbf{Lokal:} Tangente / Schmiegeparabel $\Rightarrow$ Taylor-Polynome \qquad
\textbf{Global:} Projektion / Transformation $\Rightarrow$ Fourier \& Laplace
\end{center}

Mathematik 2 zeigt damit, dass Funktionen nicht nur Rechenvorschriften sind,
sondern Elemente strukturierter R�ume, deren Eigenschaften durch Projektion,
Zerlegung und Transformation sichtbar werden.
Diese Sichtweise bildet die Grundlage f�r das mathematische Verst�ndnis
technischer Systeme.
